\scanpage{101}
\mspage{49}
% Main text, approximately lines 1--2.
In summary, we have a \deleted{\illegible{1}}\inserted{canonical} functor
% Approximately lines 3--10; the source of the vertical functor.
Isotopic category of \((X,S,\Gamma)\), \(X\) compact surface \deleted{\illegible{1}},
\(T\) oriented, \(S\) discrete subset, \(\Gamma\) operating by \unsure{\(\chi\)-automorphisms}
(\inserted{\(\sigma\in\Gamma\)}\unsure{preserving|preserves} the orientation if \(\chi(\sigma)=+1\), reversing it
if \(\chi(g)=-1\)), \(\Gamma\) transitive on \(\Pi_0(X)\),
\deleted{\illegible{1}} and if \((X_0,S_0)\) is a connected component
of \((X,S)\), we want that \deleted{\((X_0,S_0)\) \illegible{2}}\inserted{if \(X_0\simeq S^2\), we}
have \(\operatorname{card} S_{\unsure{0}}\geq 3\) \deleted{\illegible{3}}.
% Under "ait card S ≥ 3" a smaller line of text is struck through by one long stroke;
% "ait" itself is neither underlined nor struck.
\deleted{\illegible{1} \unsure{isomorphism}}
\margin{arrows: isotopy \unsure{classes} of homeomorphisms respecting all the structures.}

\diagram{\begin{tikzcd}[row sep=2.5em]
{} \arrow[d] \\
{}
\end{tikzcd}}

% Approximately lines 11--18.
Category of outer groups with \inserted{\(T\)-}loops \(\pi\),
\inserted{\(\pi\) noncommutative (i.e. \(\pi\neq 0,\mathbb Z\));}
on which a subgroup \(\Sigma\subset\Gamma\) of \(\Gamma\) \inserted{operates}
(this description being
equivalent: a description in terms
of systems of groupoids
\[
\Pi_{\unsure{D^*},\Gamma}\longrightarrow\Pi_{U,\Gamma}
\longrightarrow\Pi_{e,\Gamma}\quad\text{and }\kappa\ldots,
\]
more conceptual in certain contexts\ldots).
% Long oblique note in the left margin beside the target category.
\margin{Beware \unsure{of the} case \(S=\varnothing\) in the definition of the structures: \unsure{the} \illegible{1} (datum of
\(H^2(\pi,\mathbb Z)\overset{\kappa}{\simeq}T\)); we must of course take into account \unsure{during} the action of \(\Sigma\) on
\(H^2(\pi,\mathbb Z)\) \deleted{\unsure{and on \(T\)}}
(via its outer action on \(\pi\)) and on \(T\) (via \(\chi\)) and \unsure{require the} compatibility.}

% Approximately lines 19--25; the sentence continues on scan 102.
\deleted{\unsure{For definition}} I presume that the proof
\unsure{of the} fact that this functor is fully faithful
presents no essential difficulties,
% Oblique note in the lower-left margin, tied by a double bar and an arrow to the three preceding lines.
\margin{It would really be worth doing it very carefully\ldots}
using what is known for
\(\Gamma=1\). But the fact that, \unsure{for|\deleted{for}} \(\Gamma\) a given finite
group, \deleted{\unsure{that it is an isomorphism}}\inserted{it is \unsure{essential\textsuperscript{ly}} surjective}
% Lower-left handwritten pagination; attribution of this numbering is uncertain.
\margin{100}

\scanpage{102}
% Approximately lines 1--7; continuation from scan 101.
is an open problem on which
people are stumped. Of course we may assume \(\Sigma=\Gamma\),
\inserted{and \(\Gamma\subset\) extended Teichmüller}
and the question is whether every finite subgroup
of the extended Teichmüller is
realized as a group operating on \((X_0,S_0)\) \illegible{1}
in an essentially unique way. More
precisely, if \deleted{\illegible{1}} \(A=\) group of \deleted{\unsure{diffe}}homeomorphisms
% Approximately lines 8--16.
(or diffeomorphisms, if \(X_0\) is
differentiable) of \(X_0\), \(A^0\) its connected compo\unsure{nent|of}
(\uline{N.B.} \(A^0\) is contractible, \deleted{\unsure{)}}\inserted{in}
the anabelian cases) hence \(A/A^0=\mathcal T_{g,\nu}\)
(Teichmüller group for genus \(g\) and \(\nu\)
holes), the question comes down to: \unsure{this} whether for every
homomorphism from a finite group \(\Sigma\) into \deleted{\illegible{1}} \(\mathcal T_{g,\nu}\),
(we may assume \(\Sigma\subset\mathcal T_{g,\nu}\)), it lifts
to a homomorphism into \(A\), and whether two such
liftings are conjugate by an \(a\in A^0\)
(isotopy \inserted{in the strict sense} of two liftings \unsure{\(\sim\)}).

% Approximately lines 17--21; continuation on scan 103.
Having arrived at such a simple \deleted{description}\inserted{\unsure{reinterpretation|interpretation}}
of the 1-category (``1-isotopic'') deduced
from the 2-category of systems
\((\Pi_{D^*,\Gamma}\longrightarrow\Pi_{U,\Gamma}\longrightarrow\Pi_{e,\Gamma},\kappa)\), in terms
% Lower-left additional handwritten pagination; hand/status uncertain.
\margin{101}

\scanpage{103}
\mspage{50}
% Approximately lines 1--7; continuation from scan 102.
% The light round stroke before the square bracket may be the insertion mark for π.
of outer groups with loops (\inserted{\(\pi\)}[equipped
with a set of outer homomorphisms \(\kappa_i:T\longrightarrow\pi\), and
failing that with a \(\kappa:T\longrightarrow H^2(\pi,\mathbb Z)\))
(to do it properly, one should write
\(T^{(\otimes-1)}\simeq H^2(\pi,\mathbb Z)\), but here we have a canonical
isomorphism \(T^{\otimes-1}\simeq T\) i.e. \(T^{\otimes 2}\simeq\mathbb Z\ldots\))],
% Approximately lines 8--14.
the question arises \unsure{of how to recover},
(up to an equivalence defined up to unique isomorphism)
this diagram, in terms of \(\kappa\); everything
comes down to the \deleted{definition} description of the
categories \(\Pi_{D^*,\unsure{\Gamma}}\) and \(\Pi_{U,\Gamma}\), and the two
functors \(\Pi_{D^*,\Gamma}\longrightarrow\Pi_{U,\Gamma}\), \(\Pi_{U,\Gamma}\longrightarrow\Pi_{e,\Gamma}\);
\inserted{or what amounts to the \(=\), of the corresponding (multi-Galois) topoi and morphisms of topoi.}

% Approximately line 15.
a) \uline{Description} of \(\Pi_{U,\Gamma}\) and of \(\Pi_{U,\Gamma}\longrightarrow\Pi_{e,\Gamma}\).

% Approximately lines 16--22; continuation on scan 104.
Let more generally \(\pi\) be an \inserted{outer} group, \deleted{\unsure{such that}} \unsure{whose}
\uline{center} \uline{is trivial} (this corresponds to the
anabelian hypothesis!), let us show how we
associate with it a classifying topos
\(B_\pi\), which (as a category of sheaves) will be
% "(pour Iso)" is written under "faisceaux) sera" and tied by a curve to the place after "fonctorielle".
\(\operatorname{Ens}(\pi)\)) in a ``functorial'' way \inserted{(for \unsure{Isomorphisms|isomorphisms})}. Everything
comes down to seeing how, to a \uline{conjugacy}
% Lower-left additional handwritten pagination; hand/status uncertain.
\margin{102}

\scanpage{104}
% Approximately lines 1--7; continuation from scan 103.
\uline{class} \inserted{\(\Theta\)} of isomorphisms
\[
\pi'\xrightarrow{u}\pi,
\]
we associate an ``inverse image'' functor
\(\operatorname{Ens}(\pi)\longrightarrow\operatorname{Ens}(\pi')\), defined up to unique
isomorphism. Let us consider for every \(u\) in the class
\(\Theta\), the functor ``\(u\)-restriction of operations''
\[
\operatorname{Ens}(\pi)\xrightarrow{u^*}\operatorname{Ens}(\pi')
\]
% Approximately lines 8--11, including interlinear corrections.
% "définit," stands above the struck "est"; "avec ... ," stands above the line and is tied by a curve to the place after "donc,".
[which \deleted{\unsure{is}}\inserted{defines,} therefore, \inserted{with \(\unsure{u_*=}(u^{-1})^*\),}
an \deleted{\unsure{isomorphism}}\inserted{equivalence} of the Topoi
\[
B_{\pi'}\xrightarrow{\substack{u_\bullet\\ \sim}} B_\pi
\]
]
We shall, between these \deleted{functors}\inserted{equivalences}, \inserted{for \(u\in\Theta\),} define a
transitive system of isomorphisms (which therefore lets us
identify them with each other!).

% Approximately lines 12--20.
\unsure{Let|Let} \(u,u'\in\Theta\), whence \(u^*,u'^*\); we have by hypothesis
\deleted{definition} a \(g\in\pi\) such that \(u'=\operatorname{int}(g)\circ u\),
furthermore \(g\) is determined modulo an
element of \(\operatorname{Centr}_\pi u(\pi')=\operatorname{Centr}(\pi)=1\);
so here \(g\) is unique. But \(g\)
can be used to define a functorial
isomorphism \deleted{between \(u^*\) and \(u'^*\)}\inserted{between}
\[
i_{u',u}:u^*\xrightarrow{\sim}u'^*
\]
by taking, for \(E\in\operatorname{Ob}\operatorname{Ens}(\pi)\),
\[
i_{u',u}(E):
\underset{\substack{\unsure{\Vert}\\E}}{u^*(E)}
\longrightarrow
\underset{\substack{\unsure{\Vert}\\E}}{u'^*(E)},
\qquad i_{u',u}(E)=g_E.
\]
% Lower-left additional handwritten pagination; hand/status uncertain.
\margin{103}

\scanpage{105}
\mspage{51}
% Approximately lines 1--6.
\uline{N.B.} The argument works for every conjugacy
class of homomorphisms \inserted{of groups} \(\pi'\longrightarrow\pi\) (i.e. every outer
homomorphism) whose centralizer in \(\pi\) is
reduced to 1 (which for an epimorphism
\unsure{reduces to the} hypothesis \(\operatorname{Centr}(\pi)=1\)).
% Oblique note in the left margin. The word read "anabéliens" carries a double rule through
% or under it; deletion or emphasis is not settled.
\margin{works for \unsure{anabelian loops|anabelian groups with loops}, and the homomorphisms of such
whose image is of finite index\ldots{} (for the centralizer in such a group \(\pi\)
of a subgroup of finite \unsure{index} is still \unsure{reduced to e|reduced to 1}).}

% Approximately lines 7--14.
If now a group \(\Sigma\) operates
on the outer group \inserted{\(\pi\)}, then by the
preceding result we can say that it
also operates on the topos \(B_\pi\), \inserted{whence a topos \(B_{\pi,\Gamma}\),} (we can
also say that \deleted{\illegible{1}} Since
\(\operatorname{Centr}(\pi)=1\)) the operation of \deleted{\(\unsure{\pi}\)}\(\Sigma\) on \(\pi\)
defines an extension \(\mathcal E\) of \(\Sigma\) by \(\pi\),
% "On récupère" stands above the struck word; a curve runs from it over "bien" to the end of the line.
whence a topos \(B_{\pi,\Gamma}=B_{\mathcal E}\), \deleted{\unsure{and}} \inserted{We recover} \unsure{of}
course also \(B_{\pi,\Gamma}\longrightarrow B_{\Gamma}\). To reassure
% Approximately lines 15--22; continuation on scan 106.
ourselves, we should make sure that
if we have a homomorphism of outer groups
\(\pi'\xrightarrow{\Theta}\pi\) commuting with the action of a
\(\Sigma\), with \(\operatorname{Centr}_\pi(\Theta)=(1)\), then there exists
a \(\Theta_{\Gamma}:B_{\pi',\Gamma}\longrightarrow B_{\pi,\Gamma}\) defined up to canonical
isomorphism. Now every \(u\in\Theta\) defines a
homomorphism of extensions \inserted{\(u_{\Gamma}\) (of \(\Sigma\) by \(\pi'\) resp. \(\pi\))}
\(\mathcal E'\longrightarrow\mathcal E\) \unsure{above|below}
\(u\), and if we pass from \(u\) to \(u'=\operatorname{int}(g)\circ u\), we will have
% Lower-left additional handwritten pagination; hand/status uncertain.
\margin{104}

\scanpage{106}
% Approximately lines 1--2; continuation from scan 105.
\(u'_{\Gamma}=\operatorname{int}_{\mathcal E}(g)\circ u_{\Gamma}\), and we finish as
above, with uniqueness of \(g\).

% Approximately lines 3--8.
Now take an \inserted{\uline{injective}} outer homomorphism
\[
\deleted{\unsure{\lambda}}\Lambda=[\kappa]:T\longrightarrow\pi
\qquad (T\text{ commutative}),
\]
% Written directly under the letter Λ.
\margin{(initial of ``loops'')}
% Oblique note in the upper-left margin, tied by a curve to "injectif".
\margin{Injectivity in the case we are interested in follows from the anabelian hypothesis.}
% "supposons" stands at the start of the line, with an arrow to "que".
suppose that (for \(\kappa\in[\kappa]\))
\[
\operatorname{Centr}_\pi\kappa=\kappa(T)
\]
(which does not depend \unsure{on the|on the} choice of \(\kappa\) in
\([\kappa]\)). I will then define a
% Approximately lines 9--17.
\deleted{category} connected abelian groupoid \inserted{\(\Pi_\Lambda\)}, and an
\unsure{isomorphism} of its \(\pi_{\unsure{1}}\) with \(T\) \deleted{An object}
(hence a topos, which plays the role of
\(\widehat{\Pi}_{\unsure{D|D^*}}\)). An object will be a \(\kappa\in[\kappa]\).
A homomorphism from \(\kappa\) to \(\kappa'\) will be an
element \(g\in\operatorname{Transp}_\pi(\kappa,\kappa')\). \deleted{Let \(g\)}
The composition of homomorphisms obvious. We
do indeed get a connected groupoid, in
which \(\operatorname{Aut}(\kappa)=\operatorname{Centr}_\pi(\kappa)\overset{\inserted{\text{by hypothesis}}}{=}\kappa(T)\unsure{\xleftarrow[\sim]{\kappa}}T\).

% Approximately lines 18--21; continuation on scan 107.
% The words from "sur les topos" to "B_Λ → B_π" are squeezed between two lines.
OK. Of course, we have a homomorphism
\(\Pi_\Lambda\longrightarrow\deleted{\unsure{\pi}}\) one-point groupoid, hence
\unsure{on} the \unsure{classifying|classifying} topoi \inserted{defined by \(\pi\)} \inserted{\(B_\Lambda\longrightarrow B_\pi\)}.
We have to see the behavior of this
% Lower-left additional handwritten pagination; hand/status uncertain.
\margin{105}

\scanpage{107}
\mspage{52}
% Approximately lines 1--7; continuation from scan 106.
construction under outer homomorphisms. So let
\[
\Theta=f_\pi\deleted{\unsure{\Theta}}=[u]:\pi'\longrightarrow\pi
\]
be an outer homomorphism such that \(\operatorname{Centr}_\pi(\Theta)=1\),
and \(\Lambda':T\longrightarrow\pi'\) an injective outer homomorphism
such that \(\kappa'\in\Lambda'\Rightarrow\operatorname{Centr}_{\pi'}(\kappa')=\operatorname{Im}\kappa'\)
and let us consider the composite
\[
\deleted{\unsure{\Lambda=}}f_\pi\circ\Lambda':T\longrightarrow\pi.
\]
% Approximately lines 8--12.
Let \(d\in\mathbb Z\) and consider
\[
\Lambda=f_\pi\circ\Lambda'\circ(d.\operatorname{id}_T)
\]
\inserted{suppose that \(\kappa\in\Lambda\Rightarrow\operatorname{Centr}_\pi\kappa=\operatorname{Im}\kappa\)}
We are going to define a homomorphism from \(\Pi_{\Lambda'}\) to \(\Pi_\Lambda\),
hence a morphism of Topoi \(B_{\Lambda'}\longrightarrow B_\Lambda\),
and a commutation datum \inserted{\(\alpha\)} for the
\diagram{\begin{tikzcd}[ampersand replacement=\&,column sep=4em,row sep=3.8em]
\& B_{\Lambda'} \arrow[r,"f_D"] \arrow[d] \& B_\Lambda \arrow[d] \\
(*) \& \unsure{(B,\pi')} \arrow[r,"{\begin{gathered}\deleted{\unsure{B(f_\pi)}}\\[-2pt]\inserted{B(f_\pi)\ \text{defined above}}\end{gathered}}"']
\arrow[ur,bend left=32,shorten <=2.1em,shorten >=2.1em,"\alpha"']
\& \unsure{(B,\pi)}
\end{tikzcd}}
\margin{\illegible{2} \unsure{obviously} \illegible{5} \deleted{\unsure{of Topoi}} \illegible{2}.}
\deleted{\illegible{2} \unsure{to the} hypotheses \illegible{1} on \(\Lambda',\Theta,d\)}

% Approximately lines 13--19; continuation on scan 108.
\deleted{Let} \(\kappa':T\longrightarrow\pi'\) an object of \(\Pi_{\Lambda'}\),
we want to define (up to unique isomorphism) an
object \(\kappa\) of \(\Pi_\Lambda\). For every \(u\in f_\pi=\Theta\),
% The κ in the next formula carries no prime in the ink.
we consider \(u\circ\kappa\circ(d.\operatorname{id}_T)\), there is between them
a transitive system of isomorphisms for \(u\) varying
in \(\Theta\), we can identify them with one another.
% Lower-left additional handwritten pagination; hand/status uncertain.
\margin{106}

\scanpage{108}
% Approximately lines 1--7.
The functoriality of this object with respect to
varying \(\kappa'\) is obvious: [we can say that
\(u\in\Theta\) defines a homomorphism of groupoids
\[
u_\Lambda:\Pi_\Lambda\longrightarrow\Pi_{\Lambda'},
\]
and between these
there is a transitive system of isomorphisms \unsure{\}}. In fact,
for fixed \(u\) we have a \deleted{system} commutative
diagram of homomorphisms of groupoids
\diagram{\begin{tikzcd}[ampersand replacement=\&]
\Pi_{\Lambda'} \arrow[r,"u_\Lambda"] \arrow[d] \& \Pi_{\unsure{\Lambda|\Lambda e}} \arrow[d] \\
(e,\pi') \arrow[r,"u"] \& (e,\pi)
\end{tikzcd}}
% Approximately lines 8--11.
\deleted{Hence a}
and between these diagrams there is a
transitive system of isomorphisms, hence
the diagram \((*)\) and the commutation
datum \(\alpha\).

% Approximately lines 12--19; the final note continues on scan 109.
\deleted{The link with}
When we have a \inserted{family \(\Lambda'\) (\unsure{or a set})} \deleted{\illegible{1}} of outer
homomorphisms \(\Lambda'_{i'}\ (i'\in I'):T\longrightarrow\pi'\), and a
family \(\Lambda\unsure{=}(\Lambda_i)_{i\in I}\) of outer homomorphisms \(T\longrightarrow\pi\), and some
\(\tau:I'\longrightarrow I\), \deleted{\illegible{1}}\inserted{and \((d_{i'})_{i'\in I'}\)}
\inserted{\unsure{the|some} \(d_{i'}\in\mathbb Z\), such that}
(for \(\Theta:\pi'\longrightarrow\pi\) a given outer homomorphism)
\(\forall i'\in I'\), setting \(i=\tau(i')\), we have
% The second Λ carries no prime of its own; the prime is on the index.
\[
\Lambda_i=\Lambda_{i'}\circ\Theta\circ(d_{i'}.\operatorname{id}).
\]
[\uline{N.B.} if \(I\longrightarrow\operatorname{Homext}(T,\pi)\) injective, these conditions
show that \(\tau\) is determined by \(\Theta\), and if
% Lower-left additional handwritten pagination; hand/status uncertain.
\margin{107}

\scanpage{109}
\mspage{53}
% Approximately lines 1--6; continuation from scan 108.
in addition \(\mathbb Z\) operates faithfully on \(T\), for example
if \(T\simeq\mathbb Z\), then \((d_{i'})\) is also unique \(\ldots\)]
then we construct \deleted{\unsure{B}}
% "groupoïde" stands above the struck "Produit des", "somme des" below it.
\[
\Pi_{D^*,\Lambda'}=
\deleted{\text{Product of the}}\inserted{\text{groupoid}}\ \inserted{\text{sum of the}}\ \Pi_{\Lambda'_{i'}}
\]
and the \(=\) for \(\Pi_{D^*,\Lambda}\), and we get an
\inserted{\unsure{essentially |iso}}commutative diagram of topoi
\diagram{\begin{tikzcd}[ampersand replacement=\&]
B_{D^*,\Lambda'} \arrow[r] \arrow[d] \& B_{D^*,\Lambda} \arrow[d] \\
B_{\pi'} \arrow[r] \& B_\pi
\end{tikzcd}}

% Approximately lines 7--10.
Suppose now (whew!) that
\(\Sigma\) operates on outer groups with loops \(\ldots\)
I give up --- it is obvious that
everything works fine!

% Approximately lines 11--17.
B. \uline{Return to the arithmetic case}, if we
want to describe in ``Galoisian'' terms
pairs \((X,S)\) \deleted{\unsure{connected}} anabelian \inserted{connected} over a field
\(K\) of finite type over \(\mathbb Q\). [\uline{N.B.} if we take
a \(K\) of finite type over \(\mathbb F_p\), we would have to
restrict ourselves to the fundamental groups ``prime to
\(p\)'', apart from \unsure{that|that} our developments could
be \uline{carried out} all the \(=\) \(\ldots\)]
% Lower-left additional handwritten pagination; hand/status uncertain.
\margin{108}

\scanpage{110}
% Approximately lines 1--7.
It has become clear that in terms of an
algebraic closure \(\overline K\) of \(K\), and of the
\inserted{profinite} Galois group \(\Gamma=\operatorname{Gal}(\overline K/\unsure{K|k})\), the
simplest description is in
terms of outer groups with loops, and
of outer actions of \inserted{open} subgroups \(\Sigma\subset\Gamma\)
on them. To be precise, \deleted{let} we
% Approximately lines 8--14.
choose a connected component
\(\overline X_0\) of \(\overline X\) (or what amounts to the \(=\),
\(\overline U_0\) of \(\overline U=(X\setminus S)_{\overline K}\)), let \(\Sigma\) be its
stabilizer in \(\Gamma\) (it is replaced
by a conjugate, when we change
\(\overline U_0\)). Then \(\Sigma\) operates on the scheme
\(\unsure{\overline U_0|\overline X_0}\), hence operates by outer automorphisms on its
\(\pi_1\) (considered as an outer group).
% Approximately lines 15--18.
Now on this there is a \inserted{\(T(\overline K)\)-}loop
structure, \deleted{defined by the}\inserted{\unsure{with} as} index set
\diagram{\begin{tikzcd}[ampersand replacement=\&,row sep=0.15em,column sep=0.4em]
I=S_0(\overline K)\subset \& S(\overline K) \& \qquad[\text{empty \unsure{if|iff} }S=\varnothing] \\
\& \rotatebox{90}{$\simeq$} \& \\
\& S_0(\overline K)\wedge_\Sigma\Gamma \&
\end{tikzcd}}
% Lower-left additional handwritten pagination; hand/status uncertain.
\margin{109}
